\documentclass[11pt]{article}
\usepackage[a4paper,margin=1in]{geometry}
\usepackage[T1]{fontenc}
\usepackage[utf8]{inputenc}
\usepackage[ngerman]{babel}
\usepackage{lmodern}
\usepackage{microtype}
\usepackage{amsmath,amssymb,mathtools}
\usepackage{enumitem}
\setlength{\parindent}{1.2em}
\setlength{\parskip}{0.35em}
\setlength{\emergencystretch}{4em}
\allowdisplaybreaks
\newcommand{\dd}{\,\mathrm d}
\newcommand{\ii}{\mathrm i}
\newcommand{\cosec}{\operatorname{cosec}}
\newcommand{\art}[1]{\bigskip\begin{center}\large\bfseries #1.\end{center}\nopagebreak}
\newcommand{\tg}{\operatorname{tg}}
\newcommand{\arccosop}{\operatorname{arc\,cos}}
\begin{document}
\begin{center}
{\Large I.\\[0.4em] Über die Elemente der Theorie der Eulerschen Integrale.}\par
\medskip
{\normalsize [Inauguraldissertation, Göttingen 1852.]}
\end{center}

Es ist bekannt, daß die Ausführung der indirekten Operationen in der Analysis meist auf viel bedeutendere Schwierigkeiten stößt, als die der direkten; aber gerade dieser scheinbar unglückliche Umstand hat auf die Entwicklung der Mathematik stets den günstigsten Einfluß ausgeübt. Nicht nur, daß die Besiegung dieser Schwierigkeiten, wenn sie möglich, immer einen eigentümlichen Reiz für den Mathematiker darbietet, sondern auch gerade die Fälle, in welchen dies mit den früher eingeführten Begriffen und Hilfsmitteln nicht möglich war, haben immer der weiteren Ausbildung der Mathematik ganz neue Felder eröffnet; so führen z. B. die Operationen der Subtraktion, Division und Wurzelausziehung auf die Begriffe der negativen, gebrochenen und imaginären Zahlen, von denen jeder das Gebiet der Mathematik so außerordentlich erweitert hat. Ganz ähnlich verhält es sich nun auch in der höheren Analysis mit dem ihr zugrunde liegenden Begriff der Funktion, welcher sich anfangs nur auf die in der Elementarmathematik gelehrten Operationen (die ihnen entsprechenden Funktionen könnte man füglich Elementarfunktionen nennen) und auf deren Zusammensetzung stützt. Die Differentialrechnung, der mächtigste Hebel zur Entwicklung der Theorie der Funktionen, findet in ihrer Ausführung keine Schwierigkeiten, d. h. das Differentiieren der aus den Operationen der Elementarmathematik gebildeten Funktionen führt wieder auf eben solche Funktionen. Dagegen ist es der umgekehrten Rechnungsart, welche in ihrer Gesamtheit die Integralrechnung bildet, nur in verhältnismäßig wenigen Fällen gelungen, dasselbe zu leisten; in den meisten ist es bisher nicht geglückt, oder vielleicht auch ganz unmöglich, die Integrale gegebener Funktionen mit Hilfe eben solcher darzustellen. Aber gerade dieser Umstand hat zu einer beträchtlichen Erweiterung des Begriffs der Funktion geführt, indem man solchen nicht darstellbaren Integralen neue Namen und Bezeichnungen beigelegt, und sie dadurch in den Kreis der früheren Funktionen eingeführt hat. Bei der Entwicklung der Theorie solcher Integralfunktionen sind nun namentlich die Fälle von der größten Wichtigkeit, in denen sie sich auf die bisher allein gebräuchlichen Funktionen zurückführen lassen, indem dadurch ihr Verlauf deutlicher hervortritt, und auch oft Mittel an die Hand gegeben werden, ihre Berechnung zu erleichtern. Die Zusammenstellung dieser Fälle für die Eulerschen Integrale, mit besonderer Rücksicht auf die dabei anzuwendende Methode, ist der Hauptzweck der folgenden Abhandlung.

\art{1}

Es ist zuerst erforderlich, die Fundamentaleigenschaften der Eulerschen Integrale kurz in Erinnerung zu bringen. Die Definitionen dieser Funktionen liegen in den Gleichungen
\begin{equation}
B(a,b)=\int_0^1 x^{a-1}(1-x)^{b-1}\dd x; \qquad
\Gamma(\mu)=\int_0^\infty x^{\mu-1}e^{-x}\dd x .
\tag{1}
\end{equation}
Die Integrale rechts heißen Eulersche Integrale der ersten und der zweiten Art; die ersteren lassen sich leicht auf die letzteren zurückführen. Führt man nämlich in dem Doppelintegral
\[
\int_0^\infty\!\int_0^\infty e^{-(x+y)}x^{a-1}y^{b-1}\dd y\dd x=\Gamma(a)\Gamma(b)
\]
für $x$ und $y$ zwei neue Variabeln $r$ und $w$ ein, indem man $x+y=r$ und $x=rw$ setzt, woraus $\dd y\dd x=r\dd r\dd w$ folgt, so erhält man
\[
\int_0^\infty e^{-r}r^{a+b-1}\dd r\int_0^1 w^{a-1}(1-w)^{b-1}\dd w=\Gamma(a)\Gamma(b)
\]
und daraus zufolge der Definitionen von $B$ und $\Gamma$
\begin{equation}
B(a,b)=\frac{\Gamma(a)\Gamma(b)}{\Gamma(a+b)},
\tag{2}
\end{equation}
woraus sich zugleich ergibt, daß $B(b,a)=B(a,b)$ eine symmetrische Funktion von $a$ und $b$ ist, was man auch direkt zeigen kann, wenn man in dem Integral $B$ $(1-x)$ statt $x$ setzt.

Setzt man in dem Integral $\Gamma$ $rx$ statt $x$, und nimmt $r$ als eine positive Konstante an, so erhält man
\begin{equation}
\int_0^\infty x^{\mu-1}e^{-rx}\dd x=\frac{\Gamma(\mu)}{r^\mu}
\tag{3}
\end{equation}
und hieraus durch $n$malige partielle Differentiation in bezug auf $r$ die wichtige Relation
\begin{equation}
\Gamma(\mu+n)=(\mu+n-1)(\mu+n-2)\cdots(\mu+1)\mu\,\Gamma(\mu).
\tag{4}
\end{equation}
Es leuchtet ein, daß man $\Gamma(\mu)$ nur für jedes $\mu$ innerhalb eines Intervalls zu berechnen braucht, welches eine Einheit umfaßt, um daraus mit Hilfe dieser Gleichung $\Gamma(\mu)$ für jedes andere $\mu$ zu finden.

\art{2}

Aus den eben entwickelten Formeln lassen sich wichtige Folgerungen für die Theorie der Eulerschen Integrale ziehen, namentlich in bezug auf die Fälle, in denen sie sich ohne Hilfe neuer Funktionen darstellen lassen. Da die der ersten Art auf die der zweiten zurückgeführt werden können, so beginnen wir mit der Untersuchung der letzteren. Nun ist bekannt, daß sich bestimmte Integrale jedesmal ermitteln lassen, wenn man die unbestimmten Integrale allgemein darstellen kann (vgl. Art. 6); die Integralrechnung lehrt aber, daß dies bei dem Integral
\[
\int x^{\mu-1}e^{-x}\dd x
\]
nur dann möglich ist, wenn $\mu$ eine positive ganze Zahl $n$ ist; wir können also im voraus schließen, daß dann auch $\Gamma(n)$ sich wird angeben lassen. Nun haben wir aber in der Gleichung (4) eine Reduktionsformel gewonnen, welche uns lehrt, wie eine Gammafunktion durch eine andere dargestellt werden kann, wenn ihre Argumente um eine ganze Zahl differieren; nehmen wir also am einfachsten $\mu=1$, so erhalten wir aus der unbestimmten Integration
\begin{equation}
\Gamma(1)=\int_0^\infty e^{-x}\dd x=1 \quad\text{und folglich}\quad
\Gamma(n)=(n-1)(n-2)\cdots2\cdot1.
\tag{5}
\end{equation}
Dies ist aber auch der einzige Fall, in welchem a priori einleuchtet, daß $\Gamma(\mu)$ sich darstellen läßt.

Einen ähnlichen Weg können wir auch bei den Eulerschen Integralen der ersten Art einschlagen, indem wir zunächst die Darstellbarkeit des unbestimmten Integrals
\[
\int x^{a-1}(1-x)^{b-1}\dd x
\]
untersuchen; dieses gehört bekanntlich zu der Klasse der Integrale von sogenannten binomischen Differentialen, welche unter der allgemeinen Form
\[
\int x^m(a+bx^n)^p\dd x
\]
stehen; es ist aber bekannt, daß das binomische Differential jedesmal rational und folglich auch integrabel gemacht werden kann, wenn entweder
\[
\frac{m+1}{n}\quad\text{oder}\quad \frac{m+1}{n}+p
\]
eine ganze Zahl ist, und $m,n$ und $p$ rational sind. Damit also das obige unbestimmte Integral darstellbar sei, muß entweder $a$ [oder auch $b$, da ja $B(a,b)=B(b,a)$ ist] oder $a+b$ eine ganze Zahl sein. Der erste Fall folgt aber auch unmittelbar aus den Formeln (2) und (4); denn wenn $a$ eine ganze Zahl $m$ ist, so geben diese Formeln
\[
B(m,b)=\frac{\Gamma(m)\Gamma(b)}{\Gamma(m+b)}=
\frac{\Gamma(m)}{(m-1+b)(m-2+b)\cdots(1+b)b}
\]
und folglich mit Hilfe von der Formel (5)
\begin{equation}
B(m,b)=\frac{(m-1)(m-2)\cdots2\cdot1}{(m-1+b)(m-2+b)\cdots(1+b)b}.
\tag{6}
\end{equation}
Ebenso erhält man, wenn $n$ eine positive ganze Zahl bedeutet:
\[
B(a,n)=\frac{(n-1)(n-2)\cdots2\cdot1}{(n-1+a)(n-2+a)\cdots(1+a)a},
\]
\[
B(m,n)=\frac{(m-1)\cdots2\cdot1\,(n-1)\cdots2\cdot1}{(m+n-1)(m+n-2)\cdots2\cdot1}.
\]
Dagegen liefert der zweite Fall, in welchem $a+b$ eine ganze Zahl, ohne daß $a$ und $b$ gleichzeitig ganze Zahlen sind, unabhängig von den Eulerschen Integralen der zweiten Art, eine neue Klasse von Integralen, von denen man a priori behaupten kann, daß sie sich darstellen lassen; doch können sie alle folgendermaßen auf ein einziges zurückgeführt werden. Ist nämlich $a+b$ eine ganze positive Zahl (positive, weil als bekannt vorauszusetzen ist, daß die Eulerschen Integrale für negative Argumente stets unendlich groß ausfallen), so kann man immer $a=m+r$, $b=n-r$ setzen, worin $m$ und $n$ positive ganze Zahlen, und $r$ ein positiver echter Bruch ist. Mit Benutzung der Reduktionsformel (4) findet man dann leicht
\[
B(m+r,n-r)=\frac{\Gamma(m+r)\Gamma(n-r)}{\Gamma(m+n)}
\]
\[
=\frac{(m-1+r)(m-2+r)\cdots r\Gamma(r)\,(n-1-r)(n-2-r)\cdots(1-r)\Gamma(1-r)}{(m+n-1)(m+n-2)\cdots3\cdot2\cdot1},
\]
oder, da aus den Gleichungen (2) und (5) $\Gamma(r)\Gamma(1-r)=B(r,1-r)$ folgt,
\begin{equation}
B(m+r,n-r)=\frac{(m-1+r)\cdots r\,(n-1-r)\cdots(1-r)}{(m+n-1)(m+n-2)\cdots2\cdot1}\,B(r,1-r).
\tag{7}
\end{equation}
Das Integral $B(r,1-r)$, auf welches hierdurch die Integrale $B(m+r,n-r)$ zurückgeführt werden, ist eines der interessantesten der Integralrechnung und namentlich von der größten Wichtigkeit für die Theorie der Eulerschen Integrale beider Arten, wie schon aus der letzten Gleichung hervorgeht. So einfach aber die Form ist, in welcher der Wert desselben dargestellt werden kann, so verschiedenartig untereinander und kompliziert sind die Methoden, welche diesen Wert kennen lehren. Mit diesem Integral sollen sich daher die folgenden Artikel beschäftigen.

\art{3}

Der Gleichmäßigkeit in der Bezeichnung wegen will ich $b$ statt $r$ schreiben, und $B(b,1-b)=\Gamma(b)\Gamma(1-b)$ kurz mit $B$ bezeichnen, so daß $B$ als Funktion der einen Veränderlichen $b$ aufgefaßt wird; zwischen $b$ und $B$ besteht also folgende Gleichung:
\begin{equation}
\left\{
\begin{aligned}
B&=\int_0^1 \left(\frac{x}{1-x}\right)^b\frac{\dd x}{x}
  =\int_0^1 \left(\frac{x}{1-x}\right)^{1-b}\frac{\dd x}{x} \\
&=\int_0^\infty \frac{x^{b-1}\dd x}{x+1}
  =\int_0^\infty \frac{x^{-b}\dd x}{x+1}.
\end{aligned}
\right.
\tag{8}
\end{equation}
Die beiden letzten Formen erhält man, wenn man in den beiden ersten $x$ für $x/(1-x)$ schreibt. Wenn nun schon im vorigen Artikel $b$ auf das Intervall zwischen $0$ und $1$ beschränkt ist, weil sonst das Integral $B$ als Eulersches Integral ein negatives Argument und damit einen unendlich großen Wert erhielte, so soll in diesem Artikel die Notwendigkeit dieser Beschränkung unabhängig von der Theorie der Eulerschen Integrale in aller Strenge erwiesen werden. Größerer Allgemeinheit wegen will ich diese Untersuchung nicht unmittelbar an das Integral $B$, sondern an das allgemeinere
\[
\int_0^\infty \frac{1+x+xx+\cdots+x^{m-1}}{1+x+xx+\cdots+x^{n-1}}x^{b-1}\dd x
=\int_0^\infty \frac{1-x^m}{1-x^n}x^{b-1}\dd x
\]
anknüpfen, welches für $m=1$, $n=2$ in das Integral $B$ übergeht. Die Form rechts erfordert auch nicht einmal, daß $m$ und $n$ ganze Zahlen sind; ich will daher gleich zu der Betrachtung des Integrals
\begin{equation}
\varphi(b)=\int_0^\infty \frac{1-x^\mu}{1-x^\nu}x^{b-1}\dd x,
\tag{9}
\end{equation}
worin $b,\mu,\nu$ beliebige reelle Konstanten bedeuten sollen, übergehen.

Zuerst wird man leicht einsehen, daß man $\mu$ und $\nu$ stets positiv annehmen darf, ohne die Allgemeinheit des Integrals zu beeinträchtigen, indem man statt $1-x^\mu$ und $1-x^\nu$ auch $x^\mu(x^{-\mu}-1)$ und $x^\nu(x^{-\nu}-1)$ schreiben kann. Zerlegt man dann $\varphi(b)$ in zwei Integrale von derselben Funktion, deren Grenzen bzw. $0,1$ und $1,\infty$ sind, und setzt in dem zweiten Integral $1/x$ für $x$, so erhält auch dieses die Grenzen $0,1$, und beide lassen sich in
\begin{equation}
\varphi(b)=\int_0^1 \frac{1-x^\mu}{1-x^\nu}\left(x^b+x^{\nu-\mu-b}\right)\frac{\dd x}{x}
\tag{10}
\end{equation}
zusammenziehen. Da nun nach der Voraussetzung $\mu$ und $\nu$ positiv sind, so ist innerhalb der Integrationsgrenzen für $x$, d. h. wenn $x$ ein positiver echter Bruch ist, der Quotient
\[
\frac{1-x^\mu}{1-x^\nu}
\]
fortwährend eine positive endliche Zahl; denn auch für $x=1$ erhält dieser Quotient einen endlichen Wert, nämlich $\mu/\nu$. Bezeichnet man daher den größten und kleinsten Wert desselben mit $M$ und $N$, so sind dies ebenfalls endliche positive Zahlen, und es ist im ganzen Integrationsintervall
\[
M>\frac{1-x^\mu}{1-x^\nu}>N,
\]
und folglich auch, nach einem bekannten Satze aus der Theorie der bestimmten Integrale,
\[
M\int_0^1\left(x^b+x^{\nu-\mu-b}\right)\frac{\dd x}{x}
>\varphi(b)>
N\int_0^1\left(x^b+x^{\nu-\mu-b}\right)\frac{\dd x}{x}.
\]
Da nun das Integral, welches zu beiden Seiten von $\varphi(b)$ mit den Faktoren $M$ und $N$ vorkommt, die Werte
\[
\frac{\nu-\mu}{b(\nu-\mu-b)},\quad \pm\infty, \quad\text{oder}\quad -\infty
\]
erhält, je nachdem sowohl $b$ als $\nu-\mu-b$ positiv, oder eins von beiden $0$, oder negativ ist, so folgt, daß nur im ersten Falle $\varphi(b)$ einen endlichen, und zwar positiven, in jedem andern aber einen unendlich großen Wert erhält; es ist daher erforderlich, daß sowohl $b$ als auch $\nu-\mu-b$ eine positive Zahl sei, was man durch die Bedingung $\nu-\mu>b>0$ ausdrücken kann; hieraus geht zugleich hervor, daß $\nu>\mu$ sein muß. Wird $b=0$ oder $b=\nu-\mu$, so wird $\varphi(b)$ nicht nur unendlich groß, sondern auch unstetig, indem $\varphi(b)$ von $+\infty$ in $-\infty$ überspringt.

Aus der Gleichung (10) läßt sich noch eine interessante Folgerung ziehen; setzt man nämlich $\nu-\mu-b$ statt $b$, so erhält man unmittelbar
\begin{equation}
\varphi(\nu-\mu-b)=\varphi(b),
\tag{11}
\end{equation}
und wenn man hierin $b=(\nu-\mu)/2+b'$ setzt, in bezug auf $b'$ differentiiert und dann $b'=0$ setzt, so folgt $\varphi'((\nu-\mu)/2)=0$, worin $\varphi'(b)=\dd\varphi(b)/\dd b$ ist; um zu entscheiden, ob der Wert $b=(\nu-\mu)/2$ einem Maximum oder Minimum von $\varphi(b)$ entspricht, muß man das zweite Differentialverhältnis $\varphi''(b)$ bilden; da dieses durch das Integral
\[
\int_0^\infty \frac{1-x^\mu}{1-x^\nu}x^{b-1}\dd x\,(\ell x)^2
\]
dargestellt wird, worin $\ell x$ den natürlichen Logarithmen von $x$ bezeichnet, und folglich in dem ganzen Intervall von $b$ positiv ist, so wird $\varphi((\nu-\mu)/2)$ ein Minimum von $\varphi(b)$ sein, und zwar das einzige. Dieses Minimum wird demnach
\begin{equation}
\varphi\left(\frac{\nu}{2}-\frac{\mu}{2}\right)=
\int_0^\infty \frac{x^{\mu/2}-x^{-\mu/2}}{x^{\nu/2}-x^{-\nu/2}}\frac{\dd x}{x}
=2\int_0^\infty \frac{x^\mu-x^{-\mu}}{x^\nu-x^{-\nu}}\frac{\dd x}{x}.
\tag{12}
\end{equation}
Wenden wir das Bisherige auf unseren Fall an, in welchem $\mu=1$, $\nu=2$ zu setzen ist, so ergibt sich, daß das Integral $B$ nur dann einen endlichen, und zwar positiven Wert besitzt, wenn $b$ ein positiver echter Bruch ist; ferner daß $B=\varphi(b)=\varphi(1-b)$ ist und für $b=\tfrac12$ ein Minimum
\begin{equation}
\varphi\left(\frac12\right)=2\int_0^\infty \frac{x-1/x}{xx-1/(xx)}\frac{\dd x}{x}
=2\int_0^\infty \frac{\dd x}{xx+1}=\pi
\tag{13}
\end{equation}
erreicht.

\art{4}

Es sollen jetzt die hauptsächlichsten Beweise angegeben werden, welche bisher für den Wert des Integrals $B$ aufgestellt sind; indessen wird es genügen, kurz den Gang derselben anzudeuten, und nur da, wo eine strengere Begründung nötig scheint, näher ins Detail zu gehen.

Da schon in Art. 2 gezeigt ist, daß sich der Wert des Integrals $B(m+r,n-r)$, von welchem $B$ nur ein spezieller Fall ist, aus der unbestimmten Integration ergeben muß, wenn $r$ ein echter rationaler Bruch ist, so ist es am natürlichsten, mit dieser Methode den Anfang zu machen. Nimmt man daher $b=m/n$ an, worin $m$ und $n$ positive ganze Zahlen sind, und $m<n$, so kommt es zunächst darauf an, in dem Integral
\begin{equation}
B=\int_0^\infty \frac{x^{b-1}\dd x}{x+1}
=\int_0^\infty \frac{x^{m/n-1}\dd x}{x+1}
=n\int_0^\infty \frac{x^{m-1}\dd x}{x^n+1}
\tag{14}
\end{equation}
die unbestimmte Integration auszuführen, welche bekanntlich die Zerlegung der unter dem Integralzeichen stehenden Funktion in Partialbrüche erfordert. Setzt man zur Abkürzung $\vartheta=e^{\pi i/n}$, worin $i=\sqrt{-1}$ ist, so ist
\[
x^n+1=(x-\vartheta)(x-\vartheta^3)\cdots(x-\vartheta^{2k-1})\cdots(x-\vartheta^{2n-1});
\]
führt man danach die Zerlegung in Partialbrüche mit linearen Nennern und die Integration jedes einzelnen Gliedes aus, so findet man
\begin{equation}
n\int \frac{x^{m-1}\dd x}{x^n+1}
=-\sum_{k=1}^n \vartheta^{m(2k-1)}\ell(\vartheta^{2k-1}-x).
\tag{15}
\end{equation}
Die Summe rechts kann man auch so schreiben:
\[
-\sum_{k=1}^n \vartheta^{m(2k-1)}\ell\left(\frac{\vartheta^{2k-1}}{x}-1\right)
-\ell x\sum_{k=1}^n\vartheta^{m(2k-1)},
\]
worin
\[
\sum\vartheta^{m(2k-1)}=\vartheta^m\sum(\vartheta^{2m})^{k-1}
=\vartheta^m\frac{\vartheta^{2mn}-1}{\vartheta^{2m}-1}=0
\]
ist, so daß auch
\begin{equation}
n\int \frac{x^{m-1}\dd x}{x^n+1}
=-\sum_{k=1}^n\vartheta^{m(2k-1)}\ell\left(\frac{\vartheta^{2k-1}}{x}-1\right)
\tag{16}
\end{equation}
ist; hierbei ist wohl zu bemerken, daß die Summen in (15) und (16) vollkommen gleich, nicht etwa um eine Konstante verschieden sind. Setzt man daher in (16) $x=\infty$ und in (15) $x=0$, so gibt die Differenz das bestimmte Integral $B$, nämlich
\begin{equation}
\left\{
\begin{aligned}
B&=\sum_{k=1}^n\vartheta^{m(2k-1)}\ell(\vartheta^{2k-1})\\
&=\frac{\pi i}{n}\sum_{k=1}^n(2k-1)\vartheta^{m(2k-1)}=\frac{\pi}{\sin b\pi},
\end{aligned}
\right.
\tag{17}
\end{equation}
womit also der Wert von $B$ für rationale $b$ gefunden ist. Durch die Umformung von (15) in (16) glaube ich am kürzesten gezeigt zu haben, daß die Summe in (15) für $x=\infty$ verschwindet, und hinsichtlich seiner Strenge scheint dieser Weg denen wenigstens nicht nachzustehen, welche in den meisten Lehrbüchern der Integralrechnung befolgt sind.

\art{5}

Ein zweiter Beweis ist der folgende, welcher, so viel mir bekannt ist, von Schlömilch gegeben ist. Aus der Gleichung
\[
B=\int_0^1\frac{x^{b-1}+x^{-b}}{x+1}\dd x,
\]
welche sich aus der Formel (10) in Art. 3 ergibt, wenn man $\mu=1$, $\nu=2$ setzt, erhält man durch Entwicklung von $1/(x+1)$ nach Potenzen von $x$ mit Berücksichtigung des Restes und durch Ausführung der Integrationen für $B$ die $n$gliedrige Reihe
\begin{equation}
\frac1b+\frac{2b}{1-b^2}-\frac{2b}{4-b^2}+\frac{2b}{9-b^2}-\cdots+(-1)^n\frac{2b}{(n-1)^2-b^2}
\tag{18}
\end{equation}
nebst dem Reste
\[
\frac{(-1)^{n-1}}{n-b}+(-1)^n\int_0^1\frac{(x^{b-1}+x^{-b})x^n}{x+1}\dd x.
\]
Durch ähnliche Betrachtungen, wie die in Art. 3 über die Endlichkeit von $\varphi(b)$ angestellten, läßt sich zeigen, daß dieser Rest für unendlich wachsende $n$ gleich Null wird, so daß $B$ der unendlich fortgesetzten Reihe (18) gleich zu setzen ist. Nimmt man aber in der als bekannt vorausgesetzten Formel
\[
\cosec u=\frac1u+\frac{2u}{\pi^2-u^2}+\frac{2u}{4\pi^2-u^2}+\text{usw.}
\]
$u=b\pi$, so findet man unmittelbar für $B$ denselben Wert, wie im vorigen Artikel.

\art{6}

Ein dritter Beweis (von Cauchy) stützt sich auf einen Satz über die sogenannte Umkehrung der Integrationsordnung bei Doppelintegralen. Da die hierhergehörigen Betrachtungen sehr feiner Natur sind, und sich öfters noch Unklarheiten darüber finden, so möge es mir vergönnt sein, hier etwas weiter auszuholen, um ein sicheres Fundament für diese Untersuchung zu gewinnen. Dazu ist aber erforderlich, auf die Grundlagen der Theorie der bestimmten Integrale zurückzugehen.

Das bestimmte Integral wird meistens definiert als Differenz zweier Werte des unbestimmten Integrals, welche zwei speziellen Werten der Integrationsvariabelen entsprechen; die letztern heißen die Grenzen des Integrals. So einfach aber diese Definition scheint, so wenig kann sie strengeren Anforderungen Genüge leisten, die immer gemacht werden müssen, wenn es sich um die Festlegung einer Basis für eine ganze Theorie handelt. Der Hauptgrund für die Verwerfung dieser Definition liegt vorzüglich in dem Umstande, daß sie nicht unmittelbar auf der eigentlich gegebenen Funktion fußt, sondern als Mittelglied noch eine andere Funktion, nämlich das unbestimmte Integral voraussetzt; und dies ist ein Übelstand in mehrfacher Hinsicht. Einmal ist die Existenz des bestimmten Integrals nicht eher evident, als bis die des unbestimmten nachgewiesen ist; gesetzt aber auch, daß dies allgemein möglich wäre, so fragt sich andererseits, ob nach dieser Definition das bestimmte Integral wirklich ein bestimmtes zu nennen ist, d. h. ob es nur von der gegebenen Funktion und den Grenzen abhängt. Es ist schon mehrfach gezeigt, daß dies keineswegs der Fall ist, und es sind Fälle bekannt, in welchen diese Definition zu Zweideutigkeiten führt, welche auf diesem Wege allein gar nicht zu heben sind. Ich hoffe nun zeigen zu können, daß das nach dieser Definition aufgefaßte bestimmte Integral in jedem Falle vollkommen so unbestimmt ist wie das sogenannte unbestimmte Integral.

Während nämlich die obige Definition schon zu Zweifeln Anlaß gibt, wenn es nicht möglich ist, mit Hilfe der bekannten Methoden das unbestimmte Integral darzustellen, so geschieht dies noch in viel höherem Maße, wenn es mehrere, ja unendlich viele Funktionen gibt, deren Differential die gegebene Funktion ist. Es läßt sich zwar strenge beweisen, daß diese Funktionen nur um sogenannte Konstanten voneinander verschieden sein können, und darauf fußt gerade die obige Definition, indem sie stillschweigend voraussetzt, daß der konstante Unterschied solcher Funktionen wirklich für alle Werte der Veränderlichen derselbe bleibt. Aber gerade dies ist durchaus nicht notwendig; man kann sich hingegen denken, daß diese Konstante in verschiedenen endlichen Intervallen der Veränderlichen $x$ verschiedene Werte besitzt, und doch wird in jedem das Differentialverhältnis von $f(x)+C$ dieselbe Funktion $f'(x)$ sein, vorausgesetzt daß $C$ seinen Wert nicht stetig mit $x$ verändert, weil dann $C$ nicht mehr eine Konstante wäre. Dies leuchtet namentlich geometrisch ein, wenn man $f(x)$ als Ordinate einer krummen Linie betrachtet; man kann beliebige Stücke dieser Linie parallel der Ordinatenachse verschieben, ohne daß dadurch $f'(x)$ geändert würde. Mit andern Worten, das erste Differentialverhältnis einer Funktion gibt nicht den geringsten Aufschluß über die Stetigkeit derselben. Solche Unstetigkeiten lassen sich auch analytisch darstellen, namentlich mit Hilfe der Fourierschen Integrale, noch einfacher aber mit ganz elementaren Hilfsmitteln.

Bei der Zweideutigkeit, welche jeder Quadratwurzel hinsichtlich ihres Zeichens anhaftet, ist es durchaus erforderlich, durch ein bestimmtes Zeichen immer nur die eine Wurzel zu bezeichnen. So ist man auch darin übereingekommen, unter $\sqrt{x}$ stets die positive Quadratwurzel aus $x$ zu verstehen. Die Notwendigkeit hiervon leuchtet namentlich ein, wenn statt einer Wurzelgröße ein anderer Ausdruck, z. B. die binomische Reihe gesetzt wird, welche jedenfalls immer nur eine Wurzel ausdrückt. Dies vorausgesetzt, läßt sich leicht ein Ausdruck bilden, welcher zwar mit $x$ sich nicht stetig ändert, also eine Konstante ist, aber doch in verschiedenen Intervallen verschiedene Werte erhält. Ein solcher Ausdruck ist z. B.
\[
C\frac{x-c}{\sqrt{(x-c)^2}},
\]
welcher gleich $+C$ oder $-C$ ist, je nachdem $x$ größer oder kleiner als $c$ genommen wird. Durch Differentiation findet man natürlich
\[
\frac{\dd}{\dd x}\left(C\frac{x-c}{\sqrt{(x-c)^2}}\right)
=C\frac{\sqrt{(x-c)^2}-(x-c)\frac{x-c}{\sqrt{(x-c)^2}}}{(x-c)^2}=0
\]
und folglich ist auch
\[
\frac{\dd}{\dd x}\left(f(x)+C\frac{x-c}{\sqrt{(x-c)^2}}\right)=\frac{\dd f(x)}{\dd x}=f'(x).
\]
Wollte man daher die obige Definition des bestimmten Integrals zur Anwendung bringen, so müßte man aus dem unbestimmten Integral
\[
\int f'(x)\dd x=f(x)+C\frac{x-c}{\sqrt{(x-c)^2}}
\]
das bestimmte Integral
\[
\int_a^b f(x)\dd x=f(b)-f(a)+C\frac{b-c}{\sqrt{(b-c)^2}}-C\frac{a-c}{\sqrt{(a-c)^2}}
\]
erhalten; nimmt man hierin $a<c<b$, so ist die Summe der beiden letzten Glieder gleich $2C$, so daß das bestimmte Integral noch eine völlig willkürliche Konstante enthält.

Hiermit ist wohl das Ungenügende dieser Definition des unbestimmten Integrals dargetan, und wir wenden uns nun zu einer anderen, welche direkt von den gegebenen Größen ausgeht; nach ihr ist nämlich das bestimmte Integral als die Summe aller der unendlich kleinen Werte des gegebenen Differentials aufzufassen, wenn man der Veränderlichen $x$ stetig alle Werte beilegt, welche zwischen den gegebenen Grenzen des Integrals liegen. Diese Definition läßt wenigstens keine anderen Zweideutigkeiten zu, als solche, welche schon in der Natur der gegebenen Funktion liegen. Es läßt sich ferner zeigen, daß sie mit der ersteren stets identisch ist, sobald nur das unbestimmte Integral so gewählt ist, daß es innerhalb des Integrationsintervalls keine Unstetigkeit enthält; denn dann ist die Summe, welche nach der zweiten Definition das Integral bildet, gerade die Summe aller der unendlich kleinen Inkremente, welche das unbestimmte Integral $f(x)$ erhält, wenn $x$ das Intervall von $a$ bis $b$ stetig durchläuft. Ist aber $f(x)$ an irgend einer Stelle $c$ zwischen $a$ und $b$ unstetig, so kann man sich eine stetige Funktion $f_1(x)$ substituiert denken, deren Differential ebenfalls $f'(x)\dd x$ ist. Dann ist das bestimmte Integral $=f_1(b)-f_1(a)$, und diese Differenz unterscheidet sich von $f(b)-f(a)$ nur um den Betrag des Sprunges, welchen $f(x)$ an der Stelle $x=c$ macht, und es wird daher
\[
\int_a^b f'(x)\dd x=f(b)-f(a)+\lim\bigl(f(c-\delta)-f(c+\varepsilon)\bigr)
\]
werden, worin $\delta$ und $\varepsilon$ unendlich kleine, mit $(b-a)$ gleichstimmige Größen bedeuten; und statt dieser Gleichung kann man auch die folgende schreiben:
\[
\int_a^b f'(x)\dd x=\lim\left(\int_a^{c-\delta}f'(x)\dd x+\int_{c+\varepsilon}^b f'(x)\dd x\right).
\]

\art{7}

Die eben angestellten Betrachtungen sind vorzüglich wichtig bei solchen bestimmten Integralen, die noch eine andere Veränderliche enthalten, also bei Integralen von der Form
\[
\int_a^b f(x,\xi)\dd x,
\]
worin $\xi$ eine von $x$ unabhängige Veränderliche bedeutet. Es kann nämlich der Fall eintreten, daß das unbestimmte Integral, wenn es im allgemeinen auch für alle Werte von $x$ stetig ist, doch diese Eigenschaft verliert, wenn man der Veränderlichen $\xi$ einen bestimmten Wert beilegt. Dies ist namentlich dann zu berücksichtigen, wenn das bestimmte Integral, als Funktion von $\xi$ angesehen, einer zweiten Integration in bezug auf $\xi$ unterworfen wird, und zwar zwischen Grenzen, innerhalb deren auch der spezielle Wert von $\xi$ liegt, welcher das in bezug auf $x$ genommene unbestimmte Integral unstetig macht.

Sind $\alpha,\beta$ diese Grenzen, $c$ und $\gamma$ die Werte von $x$ und $\xi$, für welche das Integral in bezug auf $x$ unstetig wird, so muß man zufolge des vorigen Artikels
\begin{equation}
\int_\alpha^\beta \dd\xi\int_a^b f(x,\xi)\dd x=
\lim\int_\alpha^\beta\dd\xi\int_a^{c-\varepsilon} f(x,\xi)\dd x+
\int_\alpha^\beta\dd\xi\int_{c+\varepsilon}^b f(x,\xi)\dd x
\tag{19}
\end{equation}
setzen; denn man muß erst die Integration in bezug auf $x$ so ausführen, daß sie für alle Werte von $\xi$, welche bei der zweiten Integration in Betracht kommen, gültig bleibt. Ich habe diese Formel angeführt, um dadurch der unrichtigen Auffassung eines Satzes zu begegnen, der sich auf die Umkehrung der Integrationsordnung bei Doppelintegralen bezieht. In einem solchen Doppelintegral kann man nämlich die Ordnung vertauschen, also
\begin{equation}
\int_\alpha^\beta \dd\xi\int_a^b f(x,\xi)\dd x=
\int_a^b \dd x\int_\alpha^\beta f(x,\xi)\dd\xi
\tag{20}
\end{equation}
setzen, wenn $f(x,\xi)$ für alle Werte von $x$ und $\xi$ innerhalb der Integration endlich und stetig bleibt; wird aber $f(x,\xi)$ für $x=c$, $\xi=\gamma$ unstetig, so kann man noch immer
\[
\int_\alpha^\beta\dd\xi\int_a^{c-\varepsilon}f(x,\xi)\dd x
+\int_\alpha^\beta\dd\xi\int_{c+\varepsilon}^b f(x,\xi)\dd x
=\int_a^{c-\varepsilon}\dd x\int_\alpha^\beta f(x,\xi)\dd\xi
+\int_{c+\varepsilon}^b\dd x\int_\alpha^\beta f(x,\xi)\dd\xi
\]
setzen; bezeichnet man die unbestimmten Integrale in bezug auf $x$ und $\xi$ bzw. mit $F(x,\xi)$ und $\varphi(x,\xi)$ und setzt diese als stetig zwischen den Grenzen der einzelnen Integrale voraus, so geht die letzte Gleichung in
\[
\int_\alpha^\beta [F(b,\xi)-F(a,\xi)]\dd\xi-
\int_\alpha^\beta [F(c+\varepsilon,\xi)-F(c-\varepsilon,\xi)]\dd\xi
\]
\[
=\int_a^{c-\varepsilon}[\varphi(x,\beta)-\varphi(x,\alpha)]\dd x
+\int_{c+\varepsilon}^b[\varphi(x,\beta)-\varphi(x,\alpha)]\dd x
\]
über; und wenn man hierin $\varepsilon$ Null werden läßt, so erhält man
\begin{equation}
\left\{
\begin{aligned}
&\int_\alpha^\beta [F(b,\xi)-F(a,\xi)]\dd\xi
-\lim\int_\alpha^\beta [F(c+\varepsilon,\xi)-F(c-\varepsilon,\xi)]\dd\xi \\
&\hspace{7em}=\int_a^b[\varphi(x,\beta)-\varphi(x,\alpha)]\dd x .
\end{aligned}
\right.
\tag{21}
\end{equation}
Diese Formel wird meistens so aufgefaßt, als gäbe das zweite Glied auf der linken Seite den Unterschied zwischen den beiden Doppelintegralen in (20) an; aus dem im Anfang dieses Artikels Gesagten erhellt aber, daß dies nicht richtig ist, indem erst beide Glieder der linken zusammengenommen das auf der linken Seite in (20) stehende Doppelintegral darstellen. Doch wird hierdurch die Richtigkeit der Gleichung (21) nicht beeinträchtigt, und diese ist es gerade, auf welche sich der von Cauchy gegebene Beweis stützt.

Nimmt man nämlich
\[
f(x,\xi)=i f'(x+\xi i), \qquad i=\sqrt{-1}, \qquad f'(z)=\frac{\dd f(z)}{\dd z},
\]
so wird
\[
F(x,\xi)=i f(x+\xi i),\qquad \varphi(x,\xi)=f(x+\xi i)
\]
und die Gleichung (21) geht in die folgende über:
\begin{equation}
\left\{
\begin{aligned}
&i\int_\alpha^\beta [f(b+\xi i)-f(a+\xi i)]\dd\xi \\
&\quad -i\lim\int_\alpha^\beta [f(c+\varepsilon+\xi i)-f(c-\varepsilon+\xi i)]\dd\xi \\
&\hspace{8em}=\int_a^b [f(x+\beta i)-f(x+\alpha i)]\dd x .
\end{aligned}
\right.
\tag{22}
\end{equation}
Führt man in dem zweiten Gliede links eine neue Variable $\eta$ durch die Gleichung $\xi=\gamma+\varepsilon\eta$ ein, worin der Annahme nach $\alpha<\gamma<\beta$ ist, und setzt
\begin{equation}
f(z)=\frac{F(z)}{z-c-\gamma i},
\tag{23}
\end{equation}
so findet man leicht
\begin{equation}
i\lim\int_\alpha^\beta [f(c+\varepsilon+\xi i)-f(c-\varepsilon+\xi i)]\dd\xi=2\pi i F(c+\gamma i).
\tag{24}
\end{equation}

\art{8}

Aus den eben entwickelten Formeln hat nun Cauchy den Wert des Integrals $B$ abgeleitet, aber auf eine Weise, welche in einzelnen Punkten einer strengeren Begründung sehr bedürftig erscheint. Sie besteht in folgendem: Wenn die Funktion $f(z)$ so beschaffen ist, daß für jeden Wert von $\xi$ $f(\pm\infty+\xi i)=0$ und für jeden Wert von $x$ $f(x+\infty i)=0$ ist, so folgt aus den Gleichungen (22), (23) und (24), wenn man $\alpha=0$, $\beta=\infty$, $a=-\infty$, $b=\infty$ setzt,
\begin{equation}
\int_{-\infty}^{+\infty} f(x)\dd x=2\pi i F(c+\gamma i),
\tag{25}
\end{equation}
worin nun $\gamma$ zufolge der Bedingung $\alpha<\gamma<\beta$ notwendig positiv sein muß. Setzt man jetzt
\[
f(z)=\frac{(-zi)^{\mu-1}}{zz+1},
\]
worin $\mu$ eine zwischen $0$ und $2$ liegende Zahl ist (unmotiviert), so sind die Bedingungen $f(\pm\infty+\xi i)=0$ und $f(x+\infty i)=0$ erfüllt; die Werte von $x$ und $\xi$, welche $f(x+\xi i)$ unendlich machen, sind $c=0$, $\gamma=1$; es ist daher
\[
F(z)=(z-i)\frac{(-zi)^{\mu-1}}{zz+1},\qquad F(c+\gamma i)=F(i)=\frac1{2i}
\]
und folglich
\begin{equation}
\int_{-\infty}^{+\infty}\frac{(-xi)^{\mu-1}}{xx+1}\dd x=\pi.
\tag{26}
\end{equation}
Zerlegt man dies Integral in zwei andere, deren Grenzen $0,\infty$ und $-\infty,0$ sind, und setzt in dem zweiten $(-x)$ statt $x$, so findet man
\[
\int_0^\infty\frac{x^{\mu-1}\dd x}{xx+1}
=\frac{\pi}{(+i)^{\mu-1}+(-i)^{\mu-1}}
=\frac{\pi}{2\cos(\mu-1)\frac{\pi}{2}}
=\frac{\pi}{2\sin\mu\frac{\pi}{2}};
\]
und hierin braucht man bloß $x$ statt $xx$, und $\mu=2b$ (wo $b$ zwischen $0$ und $1$ liegt, wenn $0<\mu<2$ ist) zu setzen, um die Gleichung (17) wieder zu erhalten.

Hierin scheint mir namentlich die Ableitung der Gleichung (25) nicht ganz streng zu sein; denn wenn auch die Bedingungen $f(\pm\infty+\xi i)=0$, $f(x+\infty i)=0$ für jedes zwischen $0$ und $2$ liegende $\mu$ erfüllt sind (eigentlich ist dazu nur erforderlich, daß $\mu<3$ ist, so daß $\mu$ auch negativ sein kann), so ist doch bekannt, daß das Verschwinden der Funktion unter dem Integralzeichen das des Integrals nicht immer zur Folge hat, namentlich dann, wenn die eine Grenze unendlich groß ist. Ich will daher versuchen, durch die folgende Darstellung diese Zweifel zu heben und zugleich zu beweisen, daß $\mu$ zwischen den Grenzen $0$ und $2$ liegen muß.

Setzt man in der Gleichung (22) $\alpha=0$, $b=\beta=-a=k$, so geht sie mit Berücksichtigung der Gleichung (24) in folgende über
\[
i\int_0^k [f(k+\xi i)-f(-k+\xi i)]\dd\xi-2\pi i F(c+\gamma i)
=\int_{-k}^{+k} f(x+ki)\dd x-\int_{-k}^{+k} f(x)\dd x,
\]
und wenn man in dem ersten Integral $\xi=k\eta$, im zweiten $x=ky$ setzt:
\[
i\int_0^1 [f(k(1+\eta i))-f(-k(1-\eta i))]k\dd\eta-2\pi iF(c+\gamma i)
=\int_{-1}^{+1}f(k(y+i))k\dd y-\int_{-k}^{+k}f(x)\dd x.
\]
Wenn nun bei unendlichem Wachsen von $k$ die Funktionen unter den Integralzeichen verschwinden, und zwar für jeden Wert der Variabelen, so werden die Integrale selbst gleich Null. Nun ist für unseren Fall
\[
kf(k(1+\eta i))=(-i)^{\mu-1}\frac{k^\mu(1+\eta i)^{\mu-1}}{kk(1+\eta i)^2+1}
\]
und ähnlich die anderen Funktionen; damit diese Ausdrücke bei dem unendlichen Wachsen von $k$ verschwinden, ist erforderlich, daß $\mu<2$ sei, wodurch aber nicht ausgeschlossen ist, daß $\mu$ auch negativ sein kann. Jedenfalls erhält man unter dieser Annahme die Gleichung (25). Aus dem Gange des Beweises im vorigen Artikel leuchtet aber ein, daß, wenn es mehrere Paare von Werten, wie $c$ und $\gamma$ gibt, für welche $f(x,\xi)$ unendlich wird, in Gleichung (25) die Summe der ihnen entsprechenden Ausdrücke zu nehmen ist (nur mit der Bemerkung, daß, wenn $\gamma=\alpha$ ist, in Gleichung (24) $\pi iF(c+\gamma i)$ statt $2\pi iF(c+\gamma i)$ gesetzt werden muß). In unserem Falle ist aber $f(x,\xi)=i f'(x+\xi i)$ und nach der obigen Spezialisierung von $f(z)$:
\[
f'(z)=(-i)^{\mu-1}\frac{(\mu-3)z^\mu+(\mu-1)z^{\mu-2}}{(zz+1)^2},
\]
und hierin sind die komplexen Werte von $z$ aufzusuchen, welche diese Funktion unendlich machen; die ersten erhält man aus der Gleichung $zz+1=0$, woraus die beiden Systeme $(c=0,\gamma=1)$ und $(c=0,\gamma=-1)$ folgen, deren erstes oben schon behandelt ist; das zweite muß aber ausgeschlossen werden, weil dies $\gamma$ der Bedingung $\alpha<\gamma<\beta$ nicht entspricht. Wenn aber, wie eben gezeigt ist, $\mu<2$ sein muß, so ist auch $(c=0,\gamma=0)$ ein solches System, und wir hätten demnach noch die Korrektion $\pi iF(0)$ anzubringen, worin $F(z)=(z-c-\gamma i)f(z)$, also in diesem Falle
\[
F(z)=(-i)^{\mu-1}\frac{z^\mu}{zz+1}
\]
ist; für $z=0$ wird $F(z)$ nun entweder unendlich groß oder Null, je nachdem $\mu$ negativ oder positiv ist. Soll daher der Wert des Integrals in (25) endlich sein, so müssen wir das letztere annehmen, und dann ist in (26) eine Korrektion nicht mehr hinzuzufügen, da $F(0)=0$ wird. Durch diese Betrachtung ist daher die Gültigkeit der Gleichung (26), aus welcher unmittelbar der Wert des Integrals $B$ folgt, auf die Bedingung $0<\mu<2$ oder $0<b<1$ beschränkt.

\art{9}

Ein von den bisher angeführten wesentlich verschiedener Beweis ist endlich noch in der Abhandlung „Disquisitiones generales circa seriem infinitam etc. Auctore C. F. Gauß“ enthalten. Die ganze Anlage derselben macht es aber unmöglich, diesen Beweis hier darzustellen, indem die Grundlagen, auf welche er sich stützt, mit einem Schlage eine vollständige Theorie der Eulerschen Integrale ergeben, deshalb aber auch zu bedeutend sind, um hier bloß zu diesem einzigen Zwecke entwickelt zu werden.

Zu diesen will ich nun noch einen, so viel mir bekannt ist, neuen Beweis hinzufügen, der eben nicht viel Zurüstungen erfordert und sich stets in dem Gebiete der Integralrechnung hält, wenn auch die Methode nicht eben neu ist, die darauf hinaus läuft, die Aufgabe auf die Integration einer Differentialgleichung zweiter Ordnung zurückzuführen. Setzt man in der Gleichung
\[
B B=\int_0^\infty\frac{x^{b-1}\dd x}{x+1}
     \int_0^\infty\frac{y^{b-1}\dd y}{y+1}
=\int_0^\infty\frac{\dd x}{x+1}
     \int_0^\infty\frac{(xy)^{b-1}}{y+1}\dd y
\]
$y=\frac zx$, $\dd y=\frac{\dd z}{x}$, kehrt dann die Integrationsordnung um, und führt die Integration in bezug auf $x$ aus, so findet man leicht
\begin{equation}
B B=\int_0^\infty\frac{z^{b-1}\ell z}{z-1}\dd z
=\frac{\dd}{\dd b}\int_0^\infty\frac{z^{b-1}\dd z}{z-1}.
\tag{27}
\end{equation}
Durch unbestimmte Integration in bezug auf $b$ erhält man daher
\[
\int B B\,\dd b=\int_0^\infty\frac{z^{b-1}\dd z}{z-1},
\]
worin das Integral rechts sich bloß durch die Form des Nenners von dem Integral $B$ unterscheidet; und es ist leicht vorauszusehen, daß man das Integral $B$ wiedererhält, wenn man das eben gewonnene derselben Behandlung unterwirft. In der Tat findet man
\[
B\int B B\,\dd b
=\int_0^\infty\frac{\dd z}{z-1}
  \int_0^\infty\frac{(zy)^{b-1}}{y+1}\dd y,
\]
und wenn man $y=\frac xz$, $\dd y=\frac{\dd x}{z}$ setzt, dann zufolge Art. 6 und 7 das in bezug auf $z$ genommene Integral in zwei Integrale zerlegt, deren Grenzen $0,1-\delta$ und $1+\varepsilon,\infty$ sind, die Integrationsordnung umkehrt, und die Integration in bezug auf $z$ ausführt, so erhält man
\[
B\int B B\,\dd b
=\int_0^\infty\frac{x^{b-1}\ell x}{x+1}\dd x
+\lim\int_0^\infty\frac{x^{b-1}\dd x}{x+1}\,\ell\!\left(\frac{\delta}{\varepsilon}\right).
\]
Hierin ist nun zwar $\lim \ell(\delta/\varepsilon)$ unbestimmt, jedenfalls aber unabhängig von $x$ und $b$, und mag mit $k$ bezeichnet werden. Dann gibt die letzte Gleichung
\[
B\int B B\,\dd b=\frac{\dd B}{\dd b}+kB
\]
oder, wenn man bedenkt, daß in dem Integral links doch schon eine willkürliche Konstante enthalten ist,
\begin{equation}
B\int B B\,\dd b=\frac{\dd B}{\dd b},
\tag{28}
\end{equation}
woraus man durch Division mit $B$ und Differentiation in bezug auf $b$ die Differentialgleichung zweiter Ordnung
\[
B B=\frac1B\frac{\dd\dd B}{\dd b^2}
-\frac1{B B}\left(\frac{\dd B}{\dd b}\right)^2
\]
erhält, welche die Eigentümlichkeit besitzt, daß die unabhängige Variable $b$ nicht vorkommt, und sich deshalb bekanntlich auf eine Differentialgleichung erster Ordnung zurückführen läßt, wenn man das erste Differentialverhältnis als neue Variable einführt. Bezeichnen wir dieses mit $B'$, so ist
\[
\frac{\dd B}{\dd b}=B',\qquad
\frac{\dd\dd B}{\dd b^2}=\frac{\dd B'}{\dd b}
=\frac{B'\dd B'}{\dd B},
\]
und führen wir diese Transformationen in die obige Differentialgleichung ein, so geht diese in
\[
B B=\frac{B'\dd B'}{B\,\dd B}-\frac{B'B'}{B B}
\]
oder in
\[
\dd(BB)=
\frac{BB\,\dd(B'B')-B'B'\dd(BB)}{(BB)^2}
=\dd\frac{B'B'}{BB}
\]
über, deren Integral
\[
BB=cc+\frac{B'B'}{BB}
=cc+\frac1{BB}\left(\frac{\dd B}{\dd b}\right)^2
\]
ist. Zu der Bestimmung der Konstanten ist nun die Kenntnis zweier Eigenschaften des Integrals erforderlich; diese nehmen wir aus Art. 3, wo gezeigt ist, daß $B$ für $b=\frac12$ ein Minimum $=\pi$ erreicht; da gleichzeitig $\frac{\dd B}{\dd b}=0$ wird, so ergibt sich unmittelbar $cc=\pi\pi$.

Man erhält dann weiter
\[
\pm\,\dd b=
\frac{\dd B}{B\sqrt{BB-\pi\pi}}
=\frac1\pi
  \frac{-\dd\frac{\pi}{B}}
       {\sqrt{1-\frac{\pi\pi}{BB}}},
\]
folglich durch Integration
\[
\pm(b+c)=\frac1\pi\,\arccosop\frac{\pi}{B},
\qquad
B=\frac{\pi}{\cos(b+c)\pi}.
\]
Da $B$ für $b=\frac12$ den Wert $\pi$ erhält, so muß
\[
\cos\left(\frac12+c\right)\pi=-\sin c\pi=1,
\qquad \cos c\pi=0
\]
sein; folglich ist
\[
\cos(b+c)\pi=\cos b\pi\cos c\pi-\sin b\pi\sin c\pi=\sin b\pi,
\]
wodurch man wieder $B=\frac{\pi}{\sin b\pi}$ findet.

\art{10}

Weil der im vorigen Artikel gegebene Beweis den Anforderungen der größten Strenge doch noch nicht Genüge leistet, namentlich insofern das Integral
\[
\int_0^\infty\frac{x^{b-1}\dd x}{x-1}
\]
darin eine Rolle spielt, so ist es wohl nicht unangemessen, hier eine solche Modifikation noch folgen zu lassen, in welcher die unbestimmte Integration in bezug auf $b$ ganz vermieden wird. Da die Gleichung (28) sich auch so schreiben läßt
\[
\int B B\,\dd b=\frac{\dd\,\ell B}{\dd b},
\]
so läßt sich vermuten, daß eine Entwicklung des rechts stehenden Ausdrucks ebenfalls zum Ziele führt. Da $B=\Gamma(b)\Gamma(1-b)$ ist, so reicht es hin, ein Integral für $\frac{\dd\,\ell\Gamma(\mu)}{\dd\mu}$ zu finden, was sich bekanntlich auf folgendem Wege erreichen läßt. Aus der Definition von $\Gamma(\mu)$ folgt
\[
\frac{\dd\,\ell\Gamma(\mu)}{\dd\mu}
=\frac1{\Gamma(\mu)}
  \int_0^\infty x^{\mu-1}e^{-x}\ell x\,\dd x.
\]
Setzt man hierin für $\ell x$ das Integral
\[
\ell x=\int_0^\infty\frac{e^{-z}-e^{-zx}}{z}\dd z,
\]
welches sich aus dem Integral $\int_0^1y^{x-1}\dd y=\frac1x$ durch Integration in bezug auf $x$ zwischen den Grenzen $1$ und $x$, und durch Substitution von $e^{-z}$ für $y$ ergibt, so findet man durch Umkehrung der Integrationsordnung und mit Hilfe von Gleichung (3) sehr leicht
\[
\frac{\dd\,\ell\Gamma(\mu)}{\dd\mu}
=\int_0^\infty\frac{\dd z}{z}
\left(e^{-z}-\frac1{(z+1)^\mu}\right).
\]
Setzt man hierin $b$ und $(1-b)$ für $\mu$, so ergibt sich
\[
\frac{\dd\,\ell B}{\dd b}
=\int_0^\infty\frac{\dd z}{z}
\left(\frac1{(z+1)^{1-b}}-\frac1{(z+1)^b}\right),
\]
und wenn man hierin $z=x-1$ oder $z=\frac1x-1$ setzt:
\[
\frac{\dd\,\ell B}{\dd b}
=\int_1^\infty\frac{x^b-x^{1-b}}{x-1}\frac{\dd x}{x}
=\int_0^1\frac{x^b-x^{1-b}}{x-1}\frac{\dd x}{x},
\]
also auch
\[
2\frac{\dd\,\ell B}{\dd b}
=\int_0^\infty\frac{x^b-x^{1-b}}{x-1}\frac{\dd x}{x}.
\]
Vergleicht man dies mit Gleichung (27), so ergibt sich augenblicklich
\begin{equation}
2\frac{\dd\,\ell B}{\dd b}
=\int_{1-b}^{b} B B\,\dd b
=2\int_{\frac12}^{b} B B\,\dd b,
\tag{29}
\end{equation}
indem ja zufolge Art. 3 $B$ für $b=\frac12+b'$ und $b=\frac12-b'$ dieselben Werte erhält; durch Differentiation dieser Gleichung in bezug auf $b$ erhält man wieder die im vorigen Artikel behandelte Differentialgleichung.

\art{11}

Nachdem durch die angegebenen Beweise die Richtigkeit der Gleichung
\[
B(r,1-r)=\int_0^\infty\frac{x^{r-1}\dd x}{x+1}
=\frac{\pi}{\sin r\pi},
\]
worin $r$ einen positiven echten Bruch bezeichnet, außer Zweifel gesetzt ist, ergibt sich unmittelbar aus Art. 2, daß die Eulerschen Integrale der ersten Art, deren beide Argumente eine ganze positive Zahl zur Summe haben, sich wirklich darstellen lassen; man findet
\begin{equation}
B(m+r,n-r)=
\frac{(m+r-1)\cdots(1+r)r\,(n-1-r)\cdots(2-r)(1-r)}
     {(m+n-1)(m+n-2)\cdots2\cdot1}
\frac{\pi}{\sin r\pi}.
\tag{30}
\end{equation}
Man kann hiervon auch noch einen wichtigen Rückschluß auf die Eulerschen Integrale der ersten Art machen; setzt man nämlich $m=0$, $n=1$, $r=\frac12$, so findet man
\begin{equation}
\Gamma\!\left(\frac12\right)
=\int_0^\infty e^{-x}\frac{\dd x}{\sqrt{x}}
=2\int_0^\infty e^{-xx}\dd x=\sqrt{\pi}
\tag{31}
\end{equation}
und folglich nach Gleichung (4):
\begin{equation}
\Gamma\!\left(n+\frac12\right)
=(2n-1)(2n-3)\cdots5\cdot3\cdot1\,\frac{\sqrt{\pi}}{2^n}.
\tag{32}
\end{equation}

\art{12}

Nachdem in Artt. 2 und 11 die Fälle zusammengestellt sind, in welchen die Eulerschen Integrale beider Arten ohne Hilfe neuer Funktionen dargestellt werden können, will ich zum Schluß noch einmal zu dem Integral $B$ zurückkehren, um noch einige Beziehungen desselben zu anderen Integralen zu entwickeln. Unter den verschiedenen Formen, in welchen es auftritt, sind die beiden folgenden
\[
\int_{-\infty}^{+\infty}
\frac{e^{(2b-1)z}+e^{-(2b-1)z}}{e^z+e^{-z}}\dd z
\quad\text{und}\quad
2\int_0^{\frac{\pi}{2}}(\tg\varphi)^{2b-1}\dd\varphi
=2\int_0^{\frac{\pi}{2}}(\tg\varphi)^{1-2b}\dd\varphi
\]
ganz interessant; wichtiger sind aber die Verallgemeinerungen desselben durch Einführung neuer Konstanten. Dahin gehört das Integral
\begin{equation}
\int_0^\infty\frac{x^{b-1}\dd x}{x+c}
=\frac{\pi c^{b-1}}{\sin b\pi},
\tag{33}
\end{equation}
worin $c$ positiv sein muß; doch läßt sich nachträglich beweisen, daß $c$ auch imaginär sein darf; multipliziert man nämlich Zähler und Nenner der Funktion $\frac{x^{b-1}}{x+i}$ mit $(x-i)$, so zerfällt das entsprechende Integral in zwei andere, welche sich durch die Substitution $xx=y$ auf das Integral $B$ reduzieren lassen, und so findet man die Gültigkeit der Gleichung (33) für imaginäre $c$. Durch Differentiation und Integration in bezug auf $c$ kann man dann wieder eine Reihe von anderen Integralen ableiten.

Sind $c_1,c_2,\ldots,c_n$ voneinander verschiedene positive oder imaginäre Konstanten, ferner $0<b<n$ und
\[
f(x)=(x+c_1)(x+c_2)\cdots(x+c_n),\qquad
f'(x)=\frac{\dd f(x)}{\dd x},
\]
so findet man auch
\begin{equation}
\int_0^\infty\frac{x^{b-1}\dd x}{f(x)}
=\frac{\pi}{\sin b\pi}
 \sum_{k=1}^{n}\frac{c_k^{\,b-1}}{f'(-c_k)},
\tag{34}
\end{equation}
indem man $\frac{x^\beta}{f(x)}$ in Partialbrüche zerlegt, worin $\beta$ die größte in $b$ enthaltene ganze Zahl bedeutet.

Ein Vorzug des in Art. 10 gegebenen Beweises besteht auch noch darin, daß er unmittelbar eine verwandte Klasse von Integralen bestimmen lehrt. Aus den Gleichungen (27) und (29) folgt nämlich unmittelbar
\[
\int_0^\infty\frac{x^b-x^{\frac12}}{1-x}\frac{\dd x}{x}
=\pi\cot b\pi,
\]
folglich auch
\begin{equation}
\int_0^\infty\frac{x^a-x^b}{1-x}\frac{\dd x}{x}
=\pi(\cot a\pi-\cot b\pi),
\tag{35}
\end{equation}
und hieraus ergibt sich auch das in Art. 3 mit $\varphi(b)$ bezeichnete Integral
\[
\begin{aligned}
\int_0^\infty\frac{1-x^\mu}{1-x^\nu}x^{b-1}\dd x
&=\frac1\nu\int_0^\infty
  \frac{x^{\frac b\nu}-x^{\frac{\mu+b}{\nu}}}{1-x}\frac{\dd x}{x}\\
&=\frac{\pi}{\nu}
  \frac{\sin\frac{\mu\pi}{\nu}}
       {\sin\frac{b\pi}{\nu}\sin\frac{(\mu+b)\pi}{\nu}}\\
&=\frac{2\pi}{\nu}
  \frac{\sin\frac{\mu\pi}{\nu}}
       {\cos\frac{\mu\pi}{\nu}-\cos\frac{(\mu+2b)\pi}{\nu}},
\end{aligned}
\]
und das Minimum desselben
\[
\int_0^\infty
\frac{x^{\frac{\mu}{2}}-x^{-\frac{\mu}{2}}}
     {x^{\frac{\nu}{2}}-x^{-\frac{\nu}{2}}}
\frac{\dd x}{x}
=\frac{2\pi}{\nu}\,\tg\frac{\mu\pi}{2\nu}.
\]

Die meisten der hierher gehörigen Integrale finden sich in der im Auftrage des preußischen Ministeriums von Minding herausgegebenen „Sammlung von Integraltafeln“ (Berlin 1849), im Anfang der vierten Abteilung. Vielleicht ist hier der Ort, um einige Fehler, welche sich daselbst finden, anzuzeigen und zu verbessern. Durch Zerlegung von $\frac1{(x-1)(x+c)}$ in Partialbrüche und Anwendung der Formeln dieses Artikels findet man leicht
\[
\int_0^\infty\frac{x^a-1}{(x-1)(x+c)}\dd x
=\frac{\pi}{c+1}
\left(\frac{c^a-\cos a\pi}{\sin a\pi}
-\frac{\ell c}{\pi}\right),
\]
und diese Gleichung gilt für positive und negative echt gebrochene $a$, wovon man sich leicht überzeugt, wenn man $\frac1x$ und $\frac1c$ statt $x$ und $c$ schreibt. Differentiiert man in bezug auf $a$ und setzt dann $a=0$, so erhält man eine Reihe von Integralen, welche in jenen Tafeln unrichtig angegeben sind. Um die lästigen Differentiationen möglichst zu erleichtern, kann man folgenden Kunstgriff anwenden. Setzt man
\[
\frac{\ell c}{\pi}+\frac{c+1}{\pi}
\int_0^\infty\frac{x^a-1}{(x-1)(x+c)}\dd x
=\frac{c^a-\cos a\pi}{\sin a\pi}=J,
\qquad
\frac{\dd^nJ}{\dd a^n}={}^{n}J,
\]
so erhält man für $a=0$
\[
\int_0^\infty\frac{(\ell x)^n\dd x}{(x-1)(x+c)}
=\frac{\pi}{c+1}\,{}^{n}J_0.
\]
Man findet aber leicht
\[
\begin{aligned}
&J\pi^n\sin\left(a+\frac n2\right)\pi
+n\,{}^{1}J\,\pi^{n-1}\sin\left(a+\frac{n-1}{2}\right)\pi+\cdots\\
&\quad{}+n\,{}^{n-1}J\,\pi\sin\left(a+\frac12\right)\pi
+{}^{n}J\sin a\pi\\
&=\frac{\dd^n(J\sin a\pi)}{\dd a^n}
=c^a(\ell c)^n-\pi^n\cos\left(a+\frac n2\right)\pi,
\end{aligned}
\]
und für $a=0$ erhält man Rekursionsformeln für die $J_0$ mit geraden und die mit ungeraden Indizes. So findet man
\[
\begin{aligned}
\int_0^\infty\frac{\ell x\,\dd x}{(x-1)(x+c)}
&=\frac{(\ell c)^2+\pi^2}{2(c+1)},\\
\int_0^\infty\frac{(\ell x)^2\dd x}{(x-1)(x+c)}
&=\frac{\ell c\bigl((\ell c)^2+\pi^2\bigr)}{3(c+1)},\\
\int_0^\infty\frac{(\ell x)^3\dd x}{(x-1)(x+c)}
&=\frac{\bigl((\ell c)^2+\pi^2\bigr)^2}{4(c+1)},\\
\int_0^\infty\frac{(\ell x)^4\dd x}{(x-1)(x+c)}
&=\frac{\ell c\bigl((\ell c)^2+\pi^2\bigr)
        \bigl(3(\ell c)^2+7\pi^2\bigr)}{5(c+1)\cdot3},\\
\int_0^\infty\frac{(\ell x)^5\dd x}{(x-1)(x+c)}
&=\frac{\bigl((\ell c)^2+\pi^2\bigr)^2
        \bigl((\ell c)^2+3\pi^2\bigr)}{6(c+1)},
\end{aligned}
\]
während in jenen Tafeln statt der Divisoren $2,3,4,5,6$ die Divisoren $1\cdot2$, $1\cdot2\cdot3$, $1\cdot2\cdot3\cdot4$, $1\cdot2\cdot3\cdot4\cdot5$, $1\cdot2\cdot3\cdot4\cdot5\cdot6$ angegeben sind.

\end{document}
